\section{Estado del Arte}
\label{Estado del Arte}

En la actualidad, los drones han transformado profundamente la dinámica de los conflictos bélicos, como se ha evidenciado en la guerra Ruso-Ucraniana. En los primeros días del conflicto, una de las estrategias defensivas más comunes contra los drones operados remotamente por pilotos consistía en que los soldados dispararan directamente hacia ellos. Sin embargo, esta táctica demostró ser poco efectiva, con tasas de acierto extremadamente bajas debido a la velocidad y maniobrabilidad de los drones.

Para mitigar los daños ocasionados por estos dispositivos, los vehículos comenzaron a incluir mayores niveles de armadura. Este enfoque condujo al renacimiento del concepto del "tanque tortuga", el cual consiste en la adición de armadura separada del casco principal del tanque. Este diseño busca disipar las explosiones generadas por los drones, protegiendo las partes más expuestas del vehículo.

Históricamente, los tanques como el T-72, de diseño ruso, y los modelos modernos americanos, como el M1 Abrams y el M2 Bradley, fueron concebidos principalmente para enfrentamientos directos contra otros tanques, en un contexto de posible guerra entre Rusia y Estados Unidos. Alternativamente, también se diseñaron para conflictos de "contrainsurgencia", combates contra fuerzas irregulares de países con recursos limitados. No obstante, la guerra a gran escala (LSCO, por sus siglas en inglés, "Large Scale Conflict") que se está desarrollando en Ucrania ha demostrado que estos diseños enfrentan desafíos significativos frente a las tácticas modernas que emplean drones.

Una característica clave de los drones empleados en el conflicto es que operan a través del espectro radioeléctrico, utilizando frecuencias de 390-510 MHz y 750-1050 MHz. Como contramedida, se han desplegado "jammers" o bloqueadores de señal para inutilizar las comunicaciones entre los drones y sus operadores. Sin embargo, esta medida ha llevado a una carrera armamentista tecnológica, donde nuevos protocolos de comunicación se desarrollan continuamente para permitir que los drones puedan evadir estas interferencias.

En este contexto, una innovación destacable ha sido la implementación de drones rusos controlados mediante fibra óptica. Esta tecnología los hace completamente inmunes a los bloqueadores de señal, ya que elimina la dependencia del espectro radioeléctrico. Este avance representa un salto tecnológico significativo en el uso de drones en el campo de batalla, consolidando su papel como un componente esencial en las estrategias modernas de guerra.